\documentclass[10pt,a4paper]{article}

% ---------- geometry & typography ----------
\usepackage[a4paper,top=1.8cm,bottom=1.2cm,left=1.7cm,right=1.7cm,headheight=14pt,headsep=10pt]{geometry}
\usepackage[T1]{fontenc}
\usepackage[protrusion=true,expansion=false]{microtype}
\usepackage[british]{babel}

% ---------- figures, captions, lists ----------
\usepackage{graphicx}
\usepackage{caption}
\usepackage{subcaption}
\usepackage{enumitem}
\setlist[itemize]{leftmargin=1.2em,topsep=2pt,itemsep=1pt,parsep=0pt}
\setlist[enumerate]{leftmargin=1.4em,topsep=2pt,itemsep=1pt,parsep=0pt}

% ---------- section styling ----------
\usepackage{titlesec}
\titleformat{\section}{\normalfont\large\bfseries}{\thesection.}{0.5em}{}
\titlespacing*{\section}{0pt}{4pt}{1pt}

% ---------- caption styling ----------
\captionsetup{font=footnotesize,labelfont=bf,skip=2pt}
\captionsetup[sub]{font=footnotesize,labelfont=bf,skip=2pt}

% ---------- paragraph spacing ----------
\setlength{\parskip}{3pt}
\setlength{\parindent}{0pt}

% ---------- hyperref last ----------
\usepackage[hidelinks]{hyperref}

% ---------- page header ----------
\usepackage{fancyhdr}
\pagestyle{fancy}
\fancyhf{}
\fancyhead[L]{\small Antonio Galdes}
\fancyhead[R]{\small ARI5906}
\renewcommand{\headrulewidth}{0.4pt}
\setlength{\headheight}{14pt}

% ---------- title block ----------
\title{\vspace{-1.2em}\Large\bfseries Uncertainty-Aware Autonomous Control\\for Teleoperation Support}
\author{}
\date{\vspace{-3em}}

\begin{document}
\maketitle
\thispagestyle{fancy}

% ====================================================================

\section{Industry Context}
Teleoperation is being investigated at Lemonworx LTD as a way to let operators supervise deployments remotely, without being physically on site. It takes a lot of cognitive effort to maintain direct control over complicated situations, and operator tiredness reduces the amount of time an individual can remain productive. Therefore, selective automation provides a solution in which an autonomous controller handles regular manoeuvres while the operator only intervenes when necessary. Current autonomous controllers, however, do not state how confident they are in their own decisions and rarely modify their behaviour as conditions worsen. This leaves a remote operator with no clear indication for when to take over~\cite{yurtsever2020survey}.

% ====================================================================

\section{Approach}
The placement involves conducting research and creating a control policy that adjusts to the quality of its own location estimate and is aware of how confident it is in each action. A number of methods for estimating uncertainty were examined, including evidential deep learning~\cite{amini2020deep}, Bayesian dropout, and deep ensembles. Since the alternatives require multiple passes for each choice and would slow down a real-time control loop, evidential deep learning was selected since it provides a confidence estimate in a single forward pass~\cite{kendall2017uncertainties}.

% ====================================================================

\section{Preliminary Design}
The design makes use of an evidential variation of Proximal Policy Optimisation (PPO)~\cite{akgul2025eppo}, which provides a confidence estimate for each action, showing whether the system is operating under familiar conditions or is uncertain about its next course of action. Moreover, an Extended Kalman Filter (EKF) combines inertial data in the vehicle with simulated RTK-GNSS to generate location estimates. This uncertainty would then be incorporated into any policy, allowing behaviour to adapt as signal quality varies. With a four-tier signal-quality model that ranges from centimeter-accurate RTK repairs to several-meter degradation, the entire testing pipeline operates in the CARLA driving simulator~\cite{dosovitskiy2017carla}, offering a controlled testbed without committing to field trials.

% ====================================================================
\begin{figure}[h]
    \centering
    \begin{subfigure}[t]{0.48\linewidth}
        \centering
        \includegraphics[width=\linewidth,height=5cm,keepaspectratio]{rectangle}
        \caption{Layout-mode overlay rendered in windowed CARLA, showing different layout elements.\label{fig:layout}}
    \end{subfigure}
    \hfill
    \begin{subfigure}[t]{0.48\linewidth}
        \centering
        \includegraphics[width=\linewidth,height=5cm,keepaspectratio]{sensors}
        \caption{Sensor-mode side-profile annotating the IMU, RTK-GNSS antenna and 2D LiDAR mount positions.\label{fig:sensors}}
    \end{subfigure}
\end{figure}

% ====================================================================

\section{Inspector Tool}
A diagnostic inspection tool was created to assess simulation environments before any training begins, with three modes that function with any layout-and-sensor combination. Layout mode visualises feature outlines, such as parking bays and vehicle spawn sites, to identify geometric flaws before training begins (Fig.~\ref{fig:layout}). Sensor mode spawns the vehicle with each sensor mount location indicated in the simulator (Fig.~\ref{fig:sensors}), demonstrating correct physical sensor placement. Finally, live mode displays LiDAR returns in real time, highlighting the actual detection range.

% ==================================================================== 

\section{Next Steps}
The next step is designing and implementing a runtime wrapper that monitors policy confidence and sends a handover message when uncertainty exceeds a safety threshold, alerting a remote operator only when action is required. This development and validation will be carried out against synthetic confidence traces.

% ====================================================================
\clearpage
\renewcommand{\refname}{\large References}
\small
\begin{thebibliography}{9}

\bibitem{yurtsever2020survey}
E.~Yurtsever, J.~Lambert, A.~Carballo, and K.~Takeda, ``A survey of autonomous driving: Common practices and emerging technologies,'' \emph{IEEE Access}, vol.~8, pp.~58443--58469, 2020, doi: 10.1109/ACCESS.2020.2983149.

\bibitem{amini2020deep}
A.~Amini, W.~Schwarting, A.~Soleimany, and D.~Rus, ``Deep evidential regression,'' in \emph{Proc.~Conf.~Neural Inf.~Process.~Syst.~(NeurIPS)}, 2020.

\bibitem{kendall2017uncertainties}
A.~Kendall and Y.~Gal, ``What uncertainties do we need in {Bayesian} deep learning for computer vision?,'' in \emph{Proc.~Conf.~Neural Inf.~Process.~Syst.~(NeurIPS)}, 2017.

\bibitem{akgul2025eppo}
A.~Akg\"{u}l, G.~Baykal, M.~Hau{\ss}mann, and M.~Kandemir, ``Overcoming non-stationary dynamics with evidential proximal policy optimization,'' \emph{Trans.~Mach.~Learn.~Res.}, 2025, to be published.

\bibitem{dosovitskiy2017carla}
A.~Dosovitskiy, G.~Ros, F.~Codevilla, A.~L\'{o}pez, and V.~Koltun, ``{CARLA}: An open urban driving simulator,'' in \emph{Proc.~1st Annu.~Conf.~Robot Learn.~(CoRL)}, 2017.

\end{thebibliography}

\end{document}